\subsection{测试小节一}

本节测试更复杂的数学公式与范围引用.

二项式系数与分式嵌套:
\begin{equation}
    \binom{n}{k} = \frac{n!}{k!(n-k)!},
    \quad
    \frac{\frac{a}{b} + 1}{\frac{c}{d} - 1}
    \label{eq:binom-1}
\end{equation}

大运算符:
\begin{equation}
    \sum_{k=1}^{n} k = \frac{n(n+1)}{2},
    \quad
    \prod_{i=1}^{m} a_i,
    \quad
    \lim_{x \to 0} \frac{\sin x}{x} = 1,
    \quad
    \int_{-\infty}^{+\infty} e^{-x^2}\,\mathrm{d}x = \sqrt{\pi}
    \label{eq:bigop-1}
\end{equation}

split 长公式:
\begin{equation}
    \begin{split}
        (a+b+c)^3 &= a^3 + b^3 + c^3 + 3a^2b + 3a^2c \\
        &\quad + 3b^2a + 3b^2c + 3c^2a + 3c^2b + 6abc
    \end{split}
    \label{eq:split-1}
\end{equation}

连续编号公式:
\begin{equation}
    a = 1
    \label{eq:range-1}
\end{equation}

\begin{equation}
    b = 2
    \label{eq:range-2}
\end{equation}

\begin{equation}
    c = 3
    \label{eq:range-3}
\end{equation}

范围引用测试: \cref{eq:range-1,eq:range-2,eq:range-3} (连续编号应合并为范围).
